\scanpage{105}%
par un $H'$-torseur à gauche $Q$.

Moyennant la vérification laissée en suspens à l'instant, on trouve alors que la donnée d'un ``foncteur fibre \uline{filtré}'' $F$ de $\underline{C}$ sur $k$ revient à la donnée

(i) d'un foncteur fibre $F'$ \deleted{gradué de type} (d'où un \unsure{grouped|groupe} de Galois $G'=\underline{\operatorname{Aut}}(F')$ );

(ii) D'une graduation de type $\underline{Z}$ de $F'$, i.e. un homomorphisme
\[
i_1:\underline{G}_m\handwritten{\longrightarrow}G',
\]
ce qui permet de définir comme ci-dessus un sous-groupe $H'_{i_1}$ de $G'$ ;

(iii) D'un torseur à gauche $Q'$ sous $H'_{i_1}$.

A ces données, on associe simplement le foncteur fibre tordu
\[
F=F'\mathop{\handwritten{\wedge}}\limits^{H'_{i_1}}Q',
\]
$F$ étant filtré par la filtration déduite de celle de $F'$ \deleted{(\unsure{via} la \unsure{graduation|graduatio}} (associée à la graduation de $F'$) en tordant par $Q'$.

\deleted{NB Dans le cas où $F$ est le foncteur de De Rham sur la catégorie des motifs sur $k$, on prouve (moyennant le yoga général des motifs - il est possible qu'on utilise pour ceci les conjectures de Hodge ...) que le groupe $H'$ est une limite projective de groupes algébriques \deleted{\inserted{\unsure{lisses}}} \deleted{\unsure{à composantes} connexes \uline{résolubles}\unsure{.} \unsure{Donc}}\inserted{unipotents} Si on se borne à une catégorie de motifs \deleted{\unsure{de|en}} à engendrement fini, et si $G$ est connexe, on trouve aussi que $H'$ est connexe, donc étant résoluble, on trouve}

On constate aisément que le groupe $H'$ est nécessairement une limite projective de groupes algébriques \uline{unipotents}. On en conclut aussitôt que si $\underline{C}$ est à engendrement fini (ou, plus généralement, \inserted{\handwritten{à \unsure{engendrement}}} dénombrable, de façon que $G'$ \deleted{soit} donc aussi $H'$ soit limite projective d'une \uline{suite} de groupes algébriques) alors tout torseur sous $H'$ est trivial; cela signifie ici que tout foncteur fibre filtré est en fait associé à un foncteur fibre gradué (en prenant la filtration correspondant à la graduation)
% archive sheet number: 104
